% GENERATED FILE. Edit canonical lessons, then run npm run generate:books.
% canonical-source-hash: fnv1a64:46288be3221dd9c9
% canonical-lessons: TA-C405-taksi, TA-C405-metro, TA-C405-inai, TA-C405-oyvarai, TA-C405-itaivelai

\chapter{Taxis and the Metro}
\label{ch:ta-taxis-and-the-metro}
\hlchaptermodality{405}

\begin{chapteropening}
\textbf{By the end of this chapter:} \emph{I can talk about taxis, the metro, connections, a lounge and a break in a journey.}

Complete the last lesson of chapter 405: I can talk about taxis, the metro, connections, a lounge and a break in a journey.
\end{chapteropening}

\section[\texorpdfstring{\ta{டாக்ஸி} (ṭāksi)}{டாக்ஸி (ṭāksi)}]{\ta{டாக்ஸி} (ṭāksi) — a taxi}

\label{lesson:TA-C405-taksi}



\begin{quote}
Before the new one: say the Tamil for to overtake, then the Tamil for to collide.
\end{quote}

\subsection*{You'll want to know: \ta{டாக்ஸி}}
\textbf{\ta{டாக்ஸி}} — \emph{ṭāksi} — \textquotedblleft{}a taxi\textquotedblright{}.

\begin{quote}
\ta{இரவில்} \ta{நாங்கள்} \ta{டாக்ஸியில்} \ta{வீட்டுக்கு} \ta{வந்தோம்}.
\end{quote}

\begin{quote}
\emph{iravil nāṅkaḷ ṭāksiyil vīṭṭukku vantōm.}
\end{quote}

\begin{quote}
At night we came home by taxi.
\end{quote}

Before -\ta{இல்} the final -\ta{இ} takes the glide \ta{ய்}: \ta{டாக்ஸியில்}.

\begin{cousinweb}
\textbf{\ta{டாக்ஸி}} (\emph{ṭāksi}) is a loan from English \emph{taxi}.
\end{cousinweb}

\begin{grammarlens}[title={What you've built}]
One more word for getting around.
\end{grammarlens}

\subsection*{Your turn}
\begin{itemize}
\raggedright
  \item \emph{Say it:} \emph{ṭāksi}
  \item \emph{Say it:} \emph{ṭāksi}, once more
  \item \emph{Say it:} \emph{iravil nāṅkaḷ ṭāksiyil vīṭṭukku vantōm.} — At night we came home by taxi.
  \item \emph{Recall:} say the Tamil for to overtake, then the Tamil for to collide, then say \emph{ṭāksi} again
\end{itemize}

\begin{tcolorbox}[breakable,colback=teal!4,colframe=teal!35!black,title={Before you move on}]
What does \ta{டாக்ஸி} mean? (\textquotedblleft{}A taxi\textquotedblright{}.) Say it once more.
\end{tcolorbox}
\section[\texorpdfstring{\ta{மெட்ரோ} (meṭrō)}{மெட்ரோ (meṭrō)}]{\ta{மெட்ரோ} (meṭrō) — a metro; an underground railway}

\label{lesson:TA-C405-metro}



\begin{quote}
Before the new one: say the Tamil for to collide, then the Tamil for a taxi.
\end{quote}

\subsection*{You'll want to know: \ta{மெட்ரோ}}
\textbf{\ta{மெட்ரோ}} — \emph{meṭrō} — \textquotedblleft{}a metro; an underground railway\textquotedblright{}.

\begin{quote}
\ta{மெட்ரோவில்} \ta{போனால்} \ta{சீக்கிரம்} \ta{போய்ச்} \ta{சேரலாம்}.
\end{quote}

\begin{quote}
\emph{meṭrōvil pōṉāl cīkkiram pōyc cēralām.}
\end{quote}

\begin{quote}
If you go by metro, you can get there quickly.
\end{quote}

A final -\ta{ஓ} takes the glide \ta{வ்} before a case ending: \ta{மெட்ரோவில்}.

\begin{cousinweb}
\textbf{\ta{மெட்ரோ}} (\emph{meṭrō}) is a loan from English \emph{metro}.
\end{cousinweb}

\begin{grammarlens}[title={What you've built}]
One more word for getting around.
\end{grammarlens}

\subsection*{Your turn}
\begin{itemize}
\raggedright
  \item \emph{Say it:} \emph{meṭrō}
  \item \emph{Say it:} \emph{meṭrō}, once more
  \item \emph{Say it:} \emph{meṭrōvil pōṉāl cīkkiram pōyc cēralām.} — If you go by metro, you can get there quickly.
  \item \emph{Recall:} say the Tamil for to collide, then the Tamil for a taxi, then say \emph{meṭrō} again
\end{itemize}

\begin{tcolorbox}[breakable,colback=teal!4,colframe=teal!35!black,title={Before you move on}]
What does \ta{மெட்ரோ} mean? (\textquotedblleft{}A metro; an underground railway\textquotedblright{}.) Say it once more.
\end{tcolorbox}
\section[\texorpdfstring{\ta{இணை} (iṇai)}{இணை (iṇai)}]{\ta{இணை} (iṇai) — to connect; to join}

\label{lesson:TA-C405-inai}



\begin{quote}
Before the new one: say the Tamil for a taxi, then the Tamil for a metro.
\end{quote}

\subsection*{You'll want to know: \ta{இணை}}
\textbf{\ta{இணை}} — \emph{iṇai} — \textquotedblleft{}to connect; to join\textquotedblright{}.

\begin{quote}
\ta{இந்தப்} \ta{புதிய} \ta{பாலம்} \ta{இரண்டு} \ta{ஊர்களை} \ta{இணைக்கிறது}.
\end{quote}

\begin{quote}
\emph{intap putiya pālam iraṇṭu ūrkaḷai iṇaikkiṟatu.}
\end{quote}

\begin{quote}
This new bridge connects two towns.
\end{quote}

Past (I …): \ta{இணைத்தேன்}, \emph{iṇaittēṉ}.

\ta{இணைத்தேன்} is I connected something; \ta{இணைந்தேன்}, I joined.

\begin{cousinweb}
\textbf{\ta{இணை}} (\emph{iṇai}) is native Dravidian (DEDR 457); \ta{இணையம்}, \textquotedblleft{}the internet\textquotedblright{}, is built on it.
\end{cousinweb}

\begin{grammarlens}[title={What you've built}]
One more word for getting around.
\end{grammarlens}

\subsection*{Your turn}
\begin{itemize}
\raggedright
  \item \emph{Say it:} \emph{iṇai}
  \item \emph{Say it:} \emph{iṇai}, once more
  \item \emph{Say it:} \emph{intap putiya pālam iraṇṭu ūrkaḷai iṇaikkiṟatu.} — This new bridge connects two towns.
  \item \emph{Recall:} say the Tamil for a taxi, then the Tamil for a metro, then say \emph{iṇai} again
\end{itemize}

\begin{tcolorbox}[breakable,colback=teal!4,colframe=teal!35!black,title={Before you move on}]
What does \ta{இணை} mean? (\textquotedblleft{}To connect; to join\textquotedblright{}.) Say it once more.
\end{tcolorbox}
\section[\texorpdfstring{\ta{ஓய்வறை} (ōyvaṟai)}{ஓய்வறை (ōyvaṟai)}]{\ta{ஓய்வறை} (ōyvaṟai) — a lounge; a retiring room}

\label{lesson:TA-C405-oyvarai}



\begin{quote}
Before the new one: say the Tamil for a metro, then the Tamil for to connect.
\end{quote}

\subsection*{You'll want to know: \ta{ஓய்வறை}}
\textbf{\ta{ஓய்வறை}} — \emph{ōyvaṟai} — \textquotedblleft{}a lounge; a retiring room\textquotedblright{}.

\begin{quote}
\ta{பயணிகள்} \ta{ஓய்வறையில்} \ta{காத்திருந்தார்கள்}.
\end{quote}

\begin{quote}
\emph{payaṇikaḷ ōyvaṟaiyil kāttiruntārkaḷ.}
\end{quote}

\begin{quote}
The passengers waited in the lounge.
\end{quote}

\ta{ஓய்வு}, rest, + \ta{அறை}, room.

\begin{cousinweb}
\textbf{\ta{ஓய்வறை}} (\emph{ōyvaṟai}) is built from \ta{ஓய்வு}, \textquotedblleft{}rest\textquotedblright{} (DEDR 1057), and \ta{அறை}, \textquotedblleft{}a room\textquotedblright{}.
\end{cousinweb}

\begin{grammarlens}[title={What you've built}]
One more word for getting around.
\end{grammarlens}

\subsection*{Your turn}
\begin{itemize}
\raggedright
  \item \emph{Say it:} \emph{ōyvaṟai}
  \item \emph{Say it:} \emph{ōyvaṟai}, once more
  \item \emph{Say it:} \emph{payaṇikaḷ ōyvaṟaiyil kāttiruntārkaḷ.} — The passengers waited in the lounge.
  \item \emph{Recall:} say the Tamil for a metro, then the Tamil for to connect, then say \emph{ōyvaṟai} again
\end{itemize}

\begin{tcolorbox}[breakable,colback=teal!4,colframe=teal!35!black,title={Before you move on}]
What does \ta{ஓய்வறை} mean? (\textquotedblleft{}A lounge; a retiring room\textquotedblright{}.) Say it once more.
\end{tcolorbox}
\section[\texorpdfstring{\ta{இடைவேளை} (iṭaivēḷai)}{இடைவேளை (iṭaivēḷai)}]{\ta{இடைவேளை} (iṭaivēḷai) — an interval; a break}

\label{lesson:TA-C405-itaivelai}



\begin{quote}
Before the new one: say the Tamil for to connect, then the Tamil for a lounge.
\end{quote}

\subsection*{You'll want to know: \ta{இடைவேளை}}
\textbf{\ta{இடைவேளை}} — \emph{iṭaivēḷai} — \textquotedblleft{}an interval; a break\textquotedblright{}.

\begin{quote}
\ta{பயணத்தின்} \ta{நடுவில்} \ta{அரை} \ta{மணி} \ta{நேரம்} \ta{இடைவேளை} \ta{இருந்தது}.
\end{quote}

\begin{quote}
\emph{payaṇattiṉ naṭuvil arai maṇi nēram iṭaivēḷai iruntatu.}
\end{quote}

\begin{quote}
There was a half-hour break in the middle of the journey.
\end{quote}

In a cinema, \ta{இடைவேளை} is the interval.

\begin{cousinweb}
\textbf{\ta{இடைவேளை}} (\emph{iṭaivēḷai}) is built from \ta{இடை}, \textquotedblleft{}the middle, an interval\textquotedblright{} (DEDR 434), and \ta{வேளை}, \textquotedblleft{}a time of day\textquotedblright{}.
\end{cousinweb}

\begin{grammarlens}[title={What you've built}]
One more word for getting around.
\end{grammarlens}

\subsection*{Your turn}
\begin{itemize}
\raggedright
  \item \emph{Say it:} \emph{iṭaivēḷai}
  \item \emph{Say it:} \emph{iṭaivēḷai}, once more
  \item \emph{Say it:} \emph{payaṇattiṉ naṭuvil arai maṇi nēram iṭaivēḷai iruntatu.} — There was a half-hour break in the middle of the journey.
  \item \emph{Recall:} say the Tamil for to connect, then the Tamil for a lounge, then say \emph{iṭaivēḷai} again
\end{itemize}

\begin{tcolorbox}[breakable,colback=teal!4,colframe=teal!35!black,title={Before you move on}]
What does \ta{இடைவேளை} mean? (\textquotedblleft{}An interval; a break\textquotedblright{}.) Say it once more.
\end{tcolorbox}
